\chapter{Complex Amplitudes, and Where the Coherence Is}
\label{sec:quantum}
% ===================================================================

\begin{remark}[Status of this chapter]
\label{rmk:ch09status}
This chapter has been rewritten more often than any other in this part of the
book, and the reader should know what each rewriting changed.  The first
version presented the complex instance of the decoration as delivering
quantum dynamics outright.  It does not.  The second version located two
obstructions to its doing so and found one place --- the equational theory
--- where coherence could live untouched by either.  The present version
reports that the graded where-clause of Chapter~\ref{ch:grading} takes the
route the second version called ``a strictly larger construction and not a
repair of this one'', and reports what that route delivers: a construction on
which a finite, contractive, staged fragment of graded rho admits a candidate
quantum interpretation, and on which the extension to unrestricted reflective
concurrency remains open.  That is the claim, stated as the source states
it~\cite{GradedQuantum}, and it is deliberately modest.  What follows is a
construction, a boundary, and two views of where unitarity lives that this
book does not choose between.
\end{remark}

\section{The question}

Take $\Semi = \CC$.  The amplitude of a path is the product of its step
amplitudes and the transition amplitude from $P$ to $Q$ is the sum over
paths,
\[
  \langle Q \mid P \rangle \;=\; \sum_{\gamma : P \to Q} W(\gamma),
\]
with $|\langle Q \mid P\rangle|^2$ the candidate transition probability.
This much is a construction and is not in doubt.  It is also, so far, only
bookkeeping: complex numbers have been attached to steps and multiplied
along paths.

What would make it quantum is \emph{cancellation}.  Two paths $\gamma_1,
\gamma_2$ from $P$ to $Q$ with $W(\gamma_1) = -W(\gamma_2)$ should
annihilate, so that the transition has zero probability despite having two
distinct causal histories.  That is the phenomenon; everything else is
notation.  In the vocabulary of Chapter~\ref{ch:grading} it is
Principle~\ref{gw:prin:or}: a false disjunction with true disjuncts.  The
question is whether the construction can be made to have one.

\section{The first construction, and its two obstructions}
\label{sec:obstruction}

The first construction worked in the linear fragment --- weight maps keyed
by a partition --- and it is worth keeping on the page, because several other
chapters refer to it and because what it could not do is what the second
construction is for.

\begin{definition}[Configuration space, in the linear fragment]
\label{wt:def:hilbert}
Assume the set of reachable weight maps is finite.  Let
$\Hilb = \ell^2(\Cfg)$ with orthonormal basis
$\{\,|\mathfrak{c}\rangle : \mathfrak{c} \in \Cfg\,\}$ indexed by
\emph{configurations} $(P,\dec)$, not by terms.
\end{definition}

\begin{definition}[Jump operators]
\label{wt:def:jump}
For a base rule $r$ and key $\varphi \in \Phi_r$ with
$\dec(\varphi) = z = |z|e^{i\theta}$, write $\lambda(z) = |z|^2$ and put
\[
  \Lop_{r,\varphi}
  \;=\;
  e^{i\theta}\!\!\!\sum_{\mathfrak{c},\,\mathfrak{c}''}
  \sqrt{\;\sum_{(k,\sigma)\,:\,\mathfrak{c} \to \mathfrak{c}''}
        \lambda(z)\,\gfac(k)\,\chi(r,k,\sigma)\;}\;
  |\mathfrak{c}''\rangle\langle\mathfrak{c}| ,
\]
the inner sum over redexes of $r$ in class $\varphi$ carrying $\mathfrak{c}$
to $\mathfrak{c}''$.
\end{definition}

\begin{definition}[Hamiltonian from the equational theory]
\label{wt:def:hamiltonian}
Let $E_h$ be a distinguished subset of the equations of the theory,
\emph{not} quotiented into the basis, presented with real amplitudes $h_e$.
Put
\[
  H \;=\; \sum_{e \in E_h} h_e \!\!\sum_{\mathfrak{c} \equiv_e \mathfrak{c}'}\!\!
  \bigl(|\mathfrak{c}'\rangle\langle\mathfrak{c}| + |\mathfrak{c}\rangle\langle\mathfrak{c}'|\bigr),
\]
and let the \emph{quantum continuous-time Markov chain} of the weighted
theory be $(\Hilb, \Lind)$ with
$\Lind(\rho) = -i[H,\rho] + \sum_{r,\varphi}\bigl( \Lop\rho\Lop^{\dagger}
- \tfrac12\{\Lop^{\dagger}\Lop, \rho\}\bigr)$~\cite{cao2021quantum}.
\end{definition}

The square root in Definition~\ref{wt:def:jump} is outside the inner sum, and
the choice was forced by a principle the first version got right.

\begin{definition}[Conservativity]
\label{def:conservativity}
A complex weighting is \emph{conservative} over the corresponding real one
when, at zero Hamiltonian, the populations it evolves are exactly those
the real weighting evolves.
\end{definition}

\begin{proposition}[No interference between derivations; coherence only from $H$]
\label{wt:prop:interference}
Under Definition~\ref{wt:def:jump}, distinct redexes of one class with
$\equiv$-equal contracta do not interfere: their rates are summed under the
square root, so $|\langle\mathfrak{c}''|\Lop_{r,\varphi}|\mathfrak{c}\rangle|^2$
is exactly the classical rate.  Distinct classes do not interfere either,
their dissipators being summed as quantum operations.  Consequently at
$H = 0$ the populations of $(\Hilb,\Lind)$ evolve by the classical generator
of Theorem~\ref{wt:thm:markov}, and every interference effect in the model
is attributable to the commutator $-i[H,\rho]$.
\end{proposition}

\begin{theorem}[The Gillespie algorithm is the quantum jump at $H = 0$]
\label{wt:thm:degeneration}
Unraveling $(\Hilb,\Lind)$ into trajectories by the quantum-jump method
reduces at $H = 0$ to the sampling of Theorem~\ref{wt:thm:markov}, exactly.
\end{theorem}

Both results stand, and the theorem is worth having: ``Gillespie-inspired''
is a statement about a limit and not a resemblance.  But the construction
cannot do what the question asks, for two reasons that should be kept apart.

\begin{remark}[The semantic obstruction: interference versus branching time]
\label{rmk:obstruction}
The equivalence on which this book rests is bisimulation, and bisimulation
is a \emph{branching-time} notion: it distinguishes systems by the shape of
their tree of possibilities, not merely by the set of runs they admit.  Two
distinct paths to the same term are, for bisimulation, two distinct pieces
of structure, and no assignment of weights to steps can make a
branching-time equivalence identify them --- still less make them annul one
another.  Interference adds and \emph{cancels} contributions to a shared
outcome, which is a statement about runs; bisimulation refuses to pool runs
in the first place~\cite[\S13]{CE}.
\end{remark}

\begin{remark}[The constructive obstruction: conservativity kills cancellation]
\label{rmk:constructive}
Independently of the equivalence, the linear fragment attaches one amplitude
per \emph{class}.  If $m$ redexes of one class carry a configuration to one
successor, summing amplitudes and squaring would give $m^2\lambda$ where the
real instance gives $m\lambda$; conservativity forces the rates under the
root, and under the root nothing cancels.  The two paths do not cancel ---
not because branching time declines to pool them, but because the only
normalization under which the complex reading lifts the classical one adds
them where cancellation cannot occur.
\end{remark}

\begin{remark}[What the two obstructions do and do not share]
\label{rmk:two-obstructions}
They agree in their verdict and disagree in everything else.
Remark~\ref{rmk:obstruction} is a claim about \emph{observation}: a
cancellation could be computed but would not be seen.
Remark~\ref{rmk:constructive} is a claim about \emph{arithmetic}: under the
linear fragment's normalization there is no cancellation to compute.  The
second is specific to the fragment.  The first is not, and it is the one the
rest of this chapter has to answer.
\end{remark}

The second version of this chapter drew one conclusion from this pair that
survives, and one that must now be restated.  The surviving conclusion is
that a presentation has a second slot for coherence --- its equations ---
and that equations withheld from the quotient supply a Hermitian hopping
term untouched by either obstruction; that is Definition~\ref{wt:def:hamiltonian},
and Section~\ref{sec:twoviews} returns to it as one of two views.  The
conclusion to restate is that cancellation \emph{between rewrites} would need
``a decoration indexed by derivations rather than by refinement classes'', a
strictly larger construction, and that whether it could be given without
destroying locality was open.  It has been given.  It is the graded
where-clause, and locality is exactly what it costs, in a way the next
section makes precise.

\section{The second construction: amplitudes on candidates}
\label{sec:candidates}

Chapter~\ref{ch:grading} evaluates a clause once per \emph{candidate}
(Definition~\ref{gw:def:cand}) --- per occurrence-indexed match, not per
class.  Two candidates with the same contractum are two entries of the vector
$\sem{\psi}_r$, and with $\Semi = \CC$ they may carry different phases.
Equation~\eqref{gw:eq:onestep} then sums them \emph{as amplitudes}, before
anything is squared.  That is the derivation-indexed decoration the previous
version asked for, obtained not by enlarging the weight map but by discarding
the partition that made it a lookup.  Conservativity is not violated: on a
clause in normal form (Theorem~\ref{gw:thm:dnf}) every candidate of a class
carries the same value and the sum over them is the multiplicity, which is
the classical count.  Cancellation becomes available exactly when the clause
leaves normal form --- when it distinguishes candidates the partition would
have identified --- and that is the right place for it to become available,
because it is the place where the programmer has said that two ways of
firing are not the same way.

\subsection{The semantic obstruction, answered}

Remark~\ref{rmk:obstruction} said that a branching-time equivalence will not
pool runs.  It still will not, and it does not have to.
Principle~\ref{gw:prin:aperture} is the reason: the resolver is a parameter
that no context can read.  Bisimulation is a relation on terms under the
\emph{unweighted} transition relation; the Born rule is a way of
\emph{resolving} the races that relation leaves open; and the two never meet,
because the first is a statement about which races exist and the second about
how they come out.  Two paths to one normal form remain two paths for
bisimulation.  What the complex resolver adds is a distribution over normal
forms in which that normal form may have probability zero, and
Corollary~\ref{gw:cor:adequacy} says precisely what has happened: a
reachable outcome whose coefficient is a sum of nonzero terms summing to
zero.  Support adequacy holds.  Its converse fails, and the failure is the
interference.

This is the right time to say what becomes of the concession this part of
the book once carried --- that recovering interference would cost the
ontology in which bisimulation classes are the real things.  It does not,
and the reason is cleaner than the one the previous version gave.  The
ontology is a claim about what distinctions exist.  The resolver is a claim
about how the distinctions that exist get decided.  A theory may have a Born
resolver and keep every bisimulation class it had, because nothing in the
resolver identifies two terms; it only declines, with some probability, to
arrive at one of them.  Part~\ref{part:sciencebot} survives, and it survives
for a reason that does not depend on where the coherence is sited.

\subsection{What is normalized}

The route costs locality, and the cost arrives through
Definition~\ref{gw:def:contention}.  A single clause cannot be made
norm-preserving by a condition on itself, because the one-step operator sums
over sites as well as over candidates, and candidates at different sites may
be alternatives or may be concurrent.  Only alternatives are weighed against
each other; concurrent candidates both fire.  The unit of normalization is
therefore the contention component, and the alternatives within it are its
maximal matchings --- the RSpace zip.

\begin{proposition}[Completion]
\label{prop:completion}
Let $T_P$ be the one-step operator at $P$ with $\Semi = \CC$.  $T_P$ is an
isometry at $P$ if, for each contention component $C_i$,
\[
  \sum_{o \,\in\, \Out(C_i)}
  \Bigl| \sum_{\substack{m \,\text{a maximal matching of } C_i\\ \Out(m) = o}}
          w(m) \Bigr|^{2} \;=\; 1 ,
\]
together with the column-orthogonality condition of
Section~\ref{sec:fragment}.  Under staging and unit occupancy
(Definition~\ref{def:staging}) the condition is also necessary, because the
components then act on disjoint factors of the state space and $T_P$ is their
tensor product.
\end{proposition}

Three things in that statement do work, and each corrects a plausible
alternative that was measured and found wrong~\cite{GradedQuantum}.
Candidates in different components are concurrent and both fire, so treating
them as alternatives double-counts.  The alternatives within a component are
its maximal \emph{matchings}, not its individual candidates.  And the norm is
taken \emph{after} summing over the fiber of $\Out$, since matchings with
coinciding outcomes interfere before they are measured.  Normalizing over
sites, over candidates, or over matchings without the fiber sum each give the
wrong squared norm on one of five test configurations; only the fiber
convention gives $1$ on all five.

\begin{remark}[Causal histories, not interleavings]
\label{rmk:causal}
Independent redexes must be counted once.  Summing over interleavings instead
inflates the squared norm of a layer of $n$ independent sites by $n!$ ---
measured, exactly $2$, $6$ and $24$ for $n = 2,3,4$.  Under
Proposition~\ref{gw:prop:eventfact} the quotient is not imposed on
histories: concurrent firings are bundled into one matching and independent
components tensor, so the inflation cannot arise.  This is the same point
Chapter~\ref{sec:reversibility} made about the history monad --- the free
history is a tree of causal histories, not of interleavings --- arriving
from the other side.
\end{remark}

\begin{remark}[Two honest costs, restated for $\CC$]
Completion is \emph{nonlocal}: a clause can be normalized only in the light
of what else is contending, so isometry is a property of a configuration and
not of a clause.  And the sum over maximal matchings is permanent-shaped,
hence $\#P$-hard in general --- which is the physics~\cite{AaronsonArkhipov}
but bites an implementation, so what a compiler checks must be a syntactic
sufficient condition such as a bound on component size.  These are
Remark~\ref{gw:rmk:costs} with the stakes raised: in $\Rnn$ a nonlocal
normalization is an inconvenience; in $\CC$ it is the difference between a
state and a vector.
\end{remark}

\begin{conjecture}[Global norm preservation]
\label{conj:global}
Proposition~\ref{prop:completion} holds at a configuration.  That the
induced operator preserves norm along arbitrary concurrent execution, with
components merging and splitting as data are produced and consumed, is
conjectured and not proved.
\end{conjecture}

\section{When is the denotation a quantum state?}
\label{sec:fragment}

Proposition~\ref{prop:completion} makes each column of $T$ a unit vector.
That is half of $T^{\dagger}T = \mathrm{Id}$.  The other half is that distinct
columns be orthogonal, and it fails for reasons that have nothing to do with
weights.  Both failures are syntactic, and both are visible in the term.

\begin{proposition}[Absorption]
\label{prop:absorb}
A configuration with no enabled component is a fixed point of $T$ under the
absorbing convention of Definition~\ref{gw:def:denotation}.  If another
configuration in the same superposition reduces to it, the two columns are
nested and no choice of weights makes them orthogonal.  Termination is not a
unitary phenomenon.
\end{proposition}

The repair is \emph{halt tagging}: a configuration with no enabled component
does not idle; it is tagged with the step at which it halted and thereafter
frozen.  The tag records \emph{when}, never \emph{which}, so branches halting
together share it and it writes no which-path record.

\begin{definition}[Spectators and staging]
\label{def:staging}
A token not consumed by a step is a \emph{spectator}; a spectator resting on
a channel that a firing emits to is indistinguishable from the emitted
quantum, since parallel composition is a multiset.  A term is \emph{staged}
when every receipt emits only to channels that no other receipt emits to and
that no spectator occupies --- fresh wires per stage, which is what a circuit
is.
\end{definition}

\begin{conjecture}[The quantum-safe fragment]
\label{conj:safe}
Halt tagged, staged, dual-rail and token-conserving implies pairwise
orthogonal columns.  Consequently, on such terms orthogonality is automatic
and isometry reduces to Proposition~\ref{prop:completion}, which is local to
a component.
\end{conjecture}

The evidence is an exhaustive finite search over six small systems at one to
three tokens and depth three~\cite{GradedQuantum}.  Of the severe collisions
found in unrestricted terms, almost all have a normal form on one side ---
absorption --- and every non-severe one involves a spectator; under halt
tagging and staging with the dual-rail encoding, none of either kind
remains at any token count or depth searched.  That is evidence and not a
proof, and what the search does establish is the \emph{diagnosis}: the two
ways a term leaves the fragment are absorption and spectators.  That no third
way exists is exactly what is not established.

\begin{remark}[Why the circuit model could not have shown this]
Circuit-shaped terms on fresh channels have one receipt per contention
component, so a per-site normalization coincides there with
Proposition~\ref{prop:completion}, and they are staged, so they collide
nowhere.  A worked example drawn from the circuit model --- Shor's algorithm
is carried out in full in the source, dual-rail and linearly, with no
discard --- cannot exercise this part of the theory.  The fragment is where
the construction is becoming credible; the boundary is more informative than
a restriction imposed at the outset would have been.
\end{remark}

\begin{remark}[What linearity does and does not cost]
\label{rmk:linear}
A step that uses every bound name exactly once is block-diagonal in the input
basis; this is true and its hypothesis is a property of the encoding.  Under
the dual-rail encoding every one-qubit gate and every entangling two-qubit
gate is realized with no discard, so linearity does not cost universality.
What it buys is Proposition~\ref{prop:completion}'s necessity direction.
What it does not buy is recursion: recursion by reflection runs into a known
incompatibility between quantum alternation and the standard fixed-point
semantics, and the source files this as an obligation rather than a result.
The reflective calculus is not covered.  Said plainly: the quantum reading
is credible for the fragment of rho that looks like a circuit, and the part
of rho that makes it rho is the part that is open.
\end{remark}

\section{Two views of where unitarity lives}
\label{sec:coherence}
\label{sec:twoviews}

Two accounts of the same calculus are now on the table, and this book does
not choose between them.  It states both, says what each commits to, and
notes where they can be told apart.

\paragraph{View one: reductions carry coherent amplitude.}
This is the construction of Sections~\ref{sec:candidates}
and~\ref{sec:fragment}.  The graded clause at a receipt site is a matrix
over candidates and outcomes; a firing is a coherent step; interference is
between matchings of a contention component with a common outcome;
measurement is whatever reads the support.  On this view the equational
theory is quotiented into the basis as usual and contributes nothing, and
the entire quantum content is in the where-clauses.  Its strength is that
it is built, with a worked example and a measured boundary.  Its cost is
the one just stated: unrestricted reflective concurrency is open, and what
is credible is the circuit-shaped fragment.

\paragraph{View two: unitary in the equations, measurement in the reductions.}
This is the organizing intuition of a later design note~\cite{UnitaryEq},
and it is the second version of this chapter's observation taken as a
principle rather than a slot.  A GSLT presents two kinds of relation on
terms: symmetric, cost-free equations and directed, costed rewrites.  A
weighted symmetric relation can be Hermitian; a directed one cannot.  A
communication in rho \emph{quotes}: firing substitutes $@Q$ into a
continuation, the payload becomes a name, and names are compared by
equality --- so a communication writes a record, and a step that writes a
record is the kind of step Chapter~\ref{sec:reversibility} charges for.
And cost lands where information is erased, which is Landauer's principle
read operationally.  On this view the coherent dynamics is the equational
theory, suitably augmented --- messages and names become pairs of a process
and a phase, superposition enters through a binder over phase-valued
variables, couplings are inert atoms governed by path-sum equations --- and
reductions are where coherence \emph{ends}: resolving a value, or running
received code, is a measurement; only binding a receipt is coherent.  Its
strength is that it puts unitarity where the structural facts of the
calculus already point.  Its cost is that it changes the calculus ---
paired messages, a binary drop, a value binder --- and the source records it
as a design session with its decisions and open questions, not as a
semantics with a boundary measured.

\paragraph{Where they differ, and where they agree.}
They agree that Principle~\ref{wt:prin:slogan} --- \emph{a weighted theory
is quantum exactly to the extent that its presentation withholds equations
from the quotient} --- is true of view two by construction and true of view
one only vacuously, since view one withholds nothing and puts the coherence
elsewhere.  They agree that bisimulation survives, on view one by
Principle~\ref{gw:prin:aperture} and on view two because the equations that
carry phase are structural congruence enlarged, which bisimulation already
respects.  They differ on the status of a firing: coherent step or
measurement.  That difference is testable inside the calculus.  View one
predicts that a which-path record written by a reduction destroys
interference only if the record is read; view two predicts that the
reduction \emph{is} the reading.  The source for view two works a single
example --- $HZH\lvert 0\rangle = \lvert 1\rangle$ computed by equations
alone, against a which-rail detector that destroys the interference by
reductions alone --- and the two views agree on it, because in that example
the reduction does read the record.  The example that separates them has
not been written.

\begin{principle}[Where the coherence is]
\label{wt:prin:slogan}
A weighted GSLT is quantum exactly to the extent that its presentation
withholds equations from the quotient.
\end{principle}

The principle is kept, restated, and now carries a qualification: it names
the quantum content of view two and is silent about view one, whose quantum
content is in the clauses.  A reader who wants one sentence should take
this one: \emph{coherence lives either in the where-clauses or in the
equations, the two views place it differently, and nothing in this book
decides which placement is the world's.}

\begin{remark}[What is still not answered]
\label{rmk:coherence-gaps}
Three things, and they should not be run together.  First, on view one,
Conjectures~\ref{conj:global} and~\ref{conj:safe} and the recursion
obligation are open; the fragment is finite, contractive and staged, and
rho is none of those.  Second, on view two, the calculus has been redesigned
and not yet given a semantics with a measured boundary.  Third, on either
view, the notion of equivalence appropriate to a theory with coherence is
not known: bisimulation is untouched but is also uninformative about phase,
and the correctness frame of Chapter~\ref{ch:grading} ---
separations between resolvers at a fixed term --- is a frame and not yet a
theorem.  That is the precise form in which the branching-time question
returns, and it is a better form than the one this chapter started with.
\end{remark}

\section{Simulation, in the meantime}

The stochastic instance is unobstructed, and it is what the implementation
actually runs.  The classical Gillespie algorithm~\cite{gillespie1977exact}
simulates a continuous-time Markov chain by computing exit rates from the
current state, sampling a waiting time, and selecting the next rule
proportionally to rate.  With $\Semi = \Rnn$ and clauses supplying the
rates, this is a graded theory stepped forward: the simulator advances the
term, re-evaluates the clauses against it, debits the account and records
the trace as one step of one process.  By Remark~\ref{gw:rmk:state-is-term}
the chain is over terms, the state being the term.

The complex instance, on view one, is simulated by the same loop with the
contention-set completion in place of the exit rate and the Born rule in
place of the exponential clock, on the fragment where
Conjecture~\ref{conj:safe} is believed; outside it the loop runs and the
numbers it produces are not guaranteed to be those of a quantum state, which
is the sense in which the boundary of Section~\ref{sec:fragment} is the
boundary of what may be claimed.  Theorem~\ref{wt:thm:degeneration} remains
true of the linear fragment and says that the two simulators are one
simulator at $H = 0$.

What remains withheld is the interpretation and not the computation.  One
may run either loop, and the numbers coming out will be the numbers of a
quantum process on the fragment.  Whether the system they describe is a
quantum system depends on whether the amplitudes answer to anything, and
that is a question about which of the two views is the world's --- rather
than about the algorithm.  The framework has made the question askable, has
sharpened it to a single example that has not been written, and has not
answered it.

% ===================================================================
